\subsection{Domain of validity of the effective geometric description}
  \label{sec:injectivity-domain}

  The recovery of the Schwarzschild exterior geometry establishes that, in the
  presence of a localized isotropic obstruction, bounded relational flux conservation
  selects a nonlinear stationary metric configuration.
  This metric provides an effective geometric description of the projected dynamics,
  from which standard gravitational phenomena, including perihelion precession,
  follow without additional assumptions.

  It is therefore essential to characterize the domain of validity of this effective
  geometric description.

  \subsubsection*{Injectivity as a condition of geometric observability}

    The emergence of an effective metric relies on the assumption that the projection
    $\Pi : \Omega \to O$ remains locally injective at the observational scale.
    In this regime, distinct underlying configurations correspond to distinguishable
    observables, and the relational dynamics admits a geometric description in terms of
    a smooth metric $g_{\mu\nu}$.

    This condition can be expressed by comparing the projective resolution scale
    $\ell_\chi$ to the local curvature scale of the effective geometry.
    In the infrared regime, the geometric description remains valid when curvature
    invariants are small in projective units.
    A covariant injectivity condition may therefore be written schematically as
    \begin{equation}
      \mathcal{I}(x)\,\ell_\chi^2 \ll 1,
      \label{eq:injectivity-invariant}
    \end{equation}
    where $\mathcal{I}(x)$ denotes a characteristic local curvature scale.
    This relation should be understood as an order-of-magnitude criterion,
    not as an exact spectral identity.

  \subsubsection*{Application to the Schwarzschild exterior}

    For a static spherically symmetric exterior geometry, the Ricci scalar vanishes,
    and the relevant curvature scale is provided by the Kretschner invariant
    \begin{equation}
      \mathcal{K}(r)
      \equiv
      \sqrt{R_{\mu\nu\rho\sigma}R^{\mu\nu\rho\sigma}}
      \sim
      \frac{r_s}{r^3},
      \label{eq:kretschner-scale}
    \end{equation}
    with $r_s = 2G_{\mathrm{eff}}M/c^2$.

    The injectivity condition~\eqref{eq:injectivity-invariant} therefore becomes
    \begin{equation}
      \mathcal{K}(r)\,\ell_\chi^2 \ll 1.
      \label{eq:injectivity-kretsch}
    \end{equation}

    When this bound is satisfied, the projection remains locally injective and the
    effective metric description is valid.
    Conversely, when $\mathcal{K}(r)\,\ell_\chi^2 \sim 1$, distinct underlying
    configurations become observationally indistinguishable, and the geometric
    description breaks down.

  \subsubsection*{Characteristic scale of injectivity breakdown}

    The condition $\mathcal{K}(r)\,\ell_\chi^2 \sim 1$ defines a characteristic radial
    scale
    \begin{equation}
      r_{\mathrm{NI}}
      \sim
      \left(r_s \ell_\chi^2\right)^{1/3}.
      \label{eq:rNI}
    \end{equation}

    For $r \gg r_{\mathrm{NI}}$, the effective metric description is valid.
    As $r \to r_{\mathrm{NI}}$, the projective resolution limit is reached, and
    geometric observables cease to provide a complete description of the relational
    state.

    Importantly, $r_{\mathrm{NI}}$ does not in general coincide with the Schwarzschild
    radius $r_s$.
    The breakdown of geometric observability is therefore controlled by projective
    resolution rather than by the existence of a causal horizon.

  \subsubsection*{Hierarchy of regimes and observational consequences}

    The ratio between the injectivity breakdown scale and the Schwarzschild radius is
    parametrically
    \begin{equation}
      \frac{r_{\mathrm{NI}}}{r_s}
      \sim
      \left(\frac{\ell_\chi}{r_s}\right)^{2/3}.
      \label{eq:rNI-rs}
    \end{equation}

    For any microscopic value of $\ell_\chi$, one has $\ell_\chi \ll r_s$ for
    astrophysical objects, implying
    \begin{equation}
      r_{\mathrm{NI}} \ll r_s.
    \end{equation}

    The breakdown of local injectivity therefore occurs deep inside the Schwarzschild
    radius, well beyond the domain accessible to external observers.
    All currently tested gravitational phenomena, including orbital dynamics,
    light deflection, and black hole imaging, lie in the regime
    \begin{equation}
      \mathcal{K}\,\ell_\chi^2 \ll 1,
    \end{equation}
    where the effective metric description is exact.

    This provides a structural explanation for the empirical success of general
    relativity in weak and moderately strong gravitational fields, without requiring
    fine tuning or additional corrections at observable scales.

  \subsubsection*{Relation to higher-order dynamical corrections}

    The breakdown of local injectivity concerns the validity of the effective geometric
    description itself.
    This should be distinguished from the higher-order corrections to the dynamics of
    excitations discussed in Appendix~B.13, which arise within the injective regime and
    do not modify the exterior metric.

    Deviations from standard relativistic predictions are therefore expected only in
    regimes where the projection ceases to be locally injective or where strong-field
    saturation effects become relevant.

  \subsubsection*{Physical interpretation of the non-injective regime}

    The breakdown of local injectivity signals a transition in the nature of observable
    quantities.
    In the injective regime, each observable configuration corresponds to a unique
    underlying relational state, and the dynamics admits a geometric description in
    terms of a smooth metric and well-defined trajectories.

    When the condition $\mathcal{K}\,\ell_\chi^2 \ll 1$ is violated, this correspondence
    fails.
    Multiple distinct configurations of the relational substrate become indistinguishable
    under the projection $\Pi$, and the notion of a pointwise geometric state loses
    operational meaning.

    In this regime, the effective metric description ceases to be complete.
    Geometric observables such as distances, trajectories, or local curvature no longer
    provide a faithful description of the underlying dynamics.
    Instead, observables correspond to equivalence classes of configurations, reflecting
    the non-injective nature of the projection.

    This transition does not correspond to a breakdown of the underlying dynamics,
    but to a limitation of the geometric description.
    The relational substrate remains well-defined, but its evolution can no longer be
    captured by a local metric structure.

    From this perspective, the emergence of geometry is intrinsically tied to the
    resolvability of relational configurations.
    Geometry is not a fundamental structure, but an effective description valid only
    within the injective regime.

    The non-injective regime therefore represents a genuinely non-geometric phase of the
    theory, in which the appropriate description is expected to be spectral or relational
    rather than metric.

    This interpretation provides a natural framework for understanding regimes where
    classical spacetime notions are expected to fail, such as the interior of compact
    objects or the early stages of cosmological evolution, without requiring singular
    behavior of the underlying dynamics.
